\documentclass[11pt]{article}
\usepackage[a4paper,margin=28mm]{geometry}
\usepackage{amsmath,amssymb,amsthm}
\usepackage{hyperref}
\newcommand{\N}{\mathbb{N}}
\newcommand{\Z}{\mathbb{Z}}
\newtheorem{theorem}{Theorem}
\newtheorem{lemma}[theorem]{Lemma}
\newtheorem{corollary}[theorem]{Corollary}
\title{Extremal Betti Corners Form an Antichain}
\author{Source-faithful proof draft for conjecture 00000000508}
\date{8 October 2026}
\begin{document}
\maketitle

The source states that the extremal positions of $\beta_{k,k+d}$ for a
monomial ideal form an antichain. We use the standard Betti-diagram
coordinates $(k,d)$: $k$ is homological degree and $d$ is total degree
minus homological degree. The order is coordinatewise. Signed $d\in\Z$
is allowed throughout.

Let $\beta:\N\times\Z\longrightarrow\N$ be any numerical table. No
finiteness, resolution, or support hypothesis is imposed. We use the
definition in Bayer--Charalambous--Popescu~\cite{BCP}: an entry at the raw
indices $(i,j)$ is extremal when $\beta_{i,j}\ne0$ and
\[
 \beta_{l,r}=0\quad\text{for every }l\ge i,\quad r\ge j+1,
 \quad r-l\ge j-i.
\]
Write $\operatorname{Ext}(\beta)$ for the set of shifted positions
$(k,d)$ at which the raw entry $\beta_{k,k+d}$ is extremal, and put
\[
 P_\beta=\{(k,d)\in\N\times\Z:\beta_{k,k+d}\ne0\}.
\]

\begin{lemma}\label{lem:maximal}
For every numerical table $\beta$ and every $p=(k,d)$,
$p\in\operatorname{Ext}(\beta)$ if and only if $p$ is a maximal
element of $P_\beta$ in the coordinatewise order.
\end{lemma}
\begin{proof}
Suppose that $p$ is extremal and $q=(l,u)\in P_\beta$ satisfies
$p\le q$. Thus $l\ge k$ and $u\ge d$. If $q\ne p$, at least one
of these inequalities is strict. Since both coordinates are integral,
$l+u\ge k+d+1$. Taking $r=l+u$, the extremality conditions hold:
$l\ge k$, $r\ge(k+d)+1$, and $r-l=u\ge d$. They force
$\beta_{l,l+u}=0$, contradicting $q\in P_\beta$. Therefore $q=p$,
so $p$ is maximal.

Conversely, suppose that $p$ is maximal in $P_\beta$. In particular,
$\beta_{k,k+d}\ne0$. Let $l,r$ satisfy the three vanishing-region
conditions for the raw index $(k,k+d)$. If $\beta_{l,r}\ne0$, then
$q=(l,r-l)\in P_\beta$ and $p\le q$. Maximality gives $q\le p$,
so $q=p$. This implies $r=k+d$, contradicting $r\ge k+d+1$.
Consequently $\beta_{l,r}=0$, as required.
\end{proof}

\begin{theorem}\label{thm:antichain}
For every numerical table $\beta:\N\times\Z\to\N$,
$\operatorname{Ext}(\beta)$ is an antichain in $\N\times\Z$ with
the coordinatewise order.
\end{theorem}
\begin{proof}
By Lemma~\ref{lem:maximal}, its elements are precisely the maximal
elements of $P_\beta$. If two such elements $p,q$ satisfy $p\le q$,
maximality of $p$ gives $q\le p$, hence $p=q$. This excludes
comparability of distinct elements.
\end{proof}

\begin{corollary}\label{cor:decrease}
If $(k,d)$ and $(l,u)$ are extremal positions and $k<l$, then $u<d$.
\end{corollary}
\begin{proof}
If instead $d\le u$, then $(k,d)\le(l,u)$, contradicting
Theorem~\ref{thm:antichain} because their homological degrees differ.
\end{proof}

The graded Betti numbers of any monomial ideal give a numerical table
of the type in Theorem~\ref{thm:antichain}. Substitution therefore gives
the source's antichain assertion with no additional algebraic premise.
The Lean proof deliberately establishes this universal numerical fact;
it does not construct a particular ideal or assume that its extremal
support is an antichain.

\paragraph{Coordinate clarification.}
The conclusion concerns $(k,d)$ for $\beta_{k,k+d}$. It does not assert
the same product-order claim for raw indices $(k,j)$. For example, the
numerical table with only $\beta_{0,2}=\beta_{1,2}=1$ nonzero has two
extremal raw entries comparable at $(0,2)\le(1,2)$. Their shifted
positions $(0,2)$ and $(1,1)$ are incomparable. This example is only
a coordinate diagnostic; it is not asserted to be the Betti table of
a monomial ideal. The source does not explicitly spell out the order,
so acceptance requires independent review of the standard shifted
coordinate interpretation.

\paragraph{Formal verification scope.}
The corresponding Lean declarations are listed individually to keep their
names legible:
\begin{itemize}
\item \texttt{TLMC508.rawExtremal\_iff\_maximal\_shiftedSupport}
\item \texttt{TLMC508.extremal\_positions\_antichain}
\item \texttt{TLMC508.shifted\_degree\_strictly\_decreases}
\end{itemize}
Local compilation and an axiom audit are worker evidence. Deterministic
package checks and source-bound semantic acceptance remain with the manager
and independent reviewer. This document is an unaccepted source draft.

\begin{thebibliography}{1}
\bibitem{BCP}
D. Bayer, H. Charalambous, and S. Popescu,
\emph{Extremal Betti Numbers and Applications to Monomial Ideals},
Manuscript, arXiv:math/9804052.
Primary definition in the introduction, p.~1 of the manuscript:
\url{https://www.math.columbia.edu/~bayer/papers/Betti_BCP99/Betti_BCP99.pdf}.
\end{thebibliography}
\end{document}
